\chapter{Wstęp}
\label{ch:wstep}

\section{Motywacja i kontekst badawczy}
\label{sec:motywacja_i_kontekst}
Współczesny sektor usług oprogramowania (SaaS, ang. \textit{Software as a Service}) przechodzi paradygmatyczną transformację w kierunku architektur natywnych dla chmury (ang. \textit{cloud-native}). Kluczowym wyzwaniem stawianym przed systemami tej klasy jest zapewnienie wysokiej dostępności (ang. \textit{High Availability}) oraz efektywności kosztowej w warunkach skrajnie zmiennego, dynamicznego obciążenia sieciowego. Tradycyjne metody statycznej alokacji zasobów obliczeniowych prowadzą do nieefektywności: zjawisko nadwymiarowania (ang. \textit{over-provisioning}) generuje nieuzasadnione koszty infrastrukturalne, podczas gdy niedowymiarowanie (ang. \textit{under-provisioning}) skutkuje degradacją jakości usług (QoS) i naruszeniem umów SLA (ang. \textit{Service Level Agreement}).

W odpowiedzi na te problemy platformy orkiestracji kontenerów, na czele z systemem Kubernetes, wprowadziły mechanizmy automatycznego skalowania poziomego (\textit{HPA, Horizontal Pod Autoscaler}). Choć technologia ta jest powszechnie deklarowana jako panaceum na problemy wydajnościowe, jej implementacja w rzeczywistych środowiskach chmurowych napotyka bariery natury deterministycznej. Skalowanie w oparciu o standardowe metryki systemowe (zużycie procesora lub pamięci RAM) wykazuje ograniczoną skuteczność w konfrontacji z asynchronicznymi profilami operacji wejścia-wyjścia (\textit{I/O-bound}) lub gwałtownymi skokami natężenia ruchu sieciowego (\textit{Flash Crowd Effect}). Niniejsza praca podejmuje próbę szczegółowej analizy tych ograniczeń w oparciu o produkcyjne środowisko chmurowe.

\section{Problem badawczy i hipoteza}
\label{sec:problem_badawczy_hipoteza}
Głównym problemem badawczym podjętym w niniejszej pracy jest określenie wpływu doboru sygnałów telemetrycznych (metryk sterujących) na efektywność i dynamikę działania algorytmów autoskalowania poziomego w architekturach mikroserwisowych typu SaaS. W szczególności badaniu podlaga podatność pętli sprzężenia zwrotnego autoskalera na opóźnienia decyzyjne oraz narzut infrastrukturalny generowany przez dynamiczną alokację zasobów w chmurze publicznej.

W oparciu o zidentyfikowany problem sformułowano następującą **hipotezę badawczą**:
\begin{quote}
\textit{Efektywność automatycznego skalowania poziomego w systemach SaaS jest ściśle uwarunkowana korelacją pomiędzy profilem zasobowym mikroserwisu a typem metryki sterującej. Wykorzystanie niestandardowych metryk aplikacyjnych (Custom Metrics) pozwala na osiągnięcie wyższej stabilności systemu i niższej stopy błędów w porównaniu do standardowych metryk zasobowych (CPU/Memory), niezależnie od charakterystyki obciążenia aplikacyjnego.}
\end{quote}

\section{Cel i zakres pracy}
\label{sec:cel_i_zakres_pracy}
Celem głównym pracy jest zaprojektowanie, implementacja oraz empiryczna ocena wydajnościowa skalowalnej architektury systemu SaaS wdrożonej w zarządzanym klastrze kontenerów. 

Realizacja celu głównego wymagała wykonania następujących zadań cząstkowych:
\begin{enumerate}
    \item Analizy teoretycznej mechanizmów orkiestracji i algorytmów skalowania obiektów w środowisku Kubernetes.
    \item Zaprojektowania i implementacji polimorficznych mikroserwisów aplikacyjnych reprezentujących czyste profile obciążeniowe (CPU-bound oraz I/O-bound).
    \item Wdrożenia kompletnego, produkcyjnego środowiska badawczego w chmurze publicznej \textit{Google Cloud Platform (GCP)} przy użyciu usługi \textit{Google Kubernetes Engine (GKE)}.
    \item Automatyzacji procesów telemetrycznych z wykorzystaniem zaawansowanego stosu monitorującego opartego na narzędziach \textit{Prometheus}, \textit{Prometheus Adapter} oraz \textit{Grafana}.
    \item Przeprowadzenia wieloetapowych testów obciążeniowych za pomocą rozproszonego narzędzia \textit{Grafana k6} w ramach zdefiniowanej macierzy 8 scenariuszy eksperymentalnych.
    \item Spisania, sparametryzowania i przeanalizowania ujawnionych anomalii wydajnościowych sieciowo-sprzętowych.
\end{enumerate}

Zakres pracy obejmuje warstwę architektoniczną systemu (FastAPI, Docker), warstwę orkiestracji i infrastruktury sieciowej VPC w chmurze GCP, systemy zbierania szeregów czasowych oraz logikę skryptów generowania ruchu HTTP. Poza zakresem pracy pozostaje analiza skalowania pionowego (VPA) oraz badanie rozproszonych baz danych (stateful).

\section{Architektury skalowalne systemów SaaS}
\label{sec:architektury_skalowalne}

\subsection{Monolit, SOA i mikroserwisy}
\label{subsec:monolit_mikroserwisy}
Ewolucja struktur architektonicznych systemów informatycznych doprowadziła do wyparcia klasycznego wzorca monolitycznego na rzecz luźno powiązanych mikroserwisów. W architekturze monolitycznej wszystkie komponenty aplikacyjne (logika biznesowa, warstwa prezentacji, dostęp do danych) współdzielą jedną przestrzeń adresową pamięci i są uruchamiane jako pojedynczy proces. Podejście to utrudnia skalowanie: zwiększenie wydajności jednego przeciążonego elementu wymaga powielenia całego monolitu, co drastycznie zwiększa narzut pamięciowy.

Architektura mikroserwisowa dekomponuje system na zbiór autonomicznych, wyspecjalizowanych usług komunikujących się za pomocą lekkich protokołów sieciowych (np. HTTP/REST, gRPC). Każdy mikroserwis posiada niezależny cykl życia, własną domenę danych oraz dedykowane limity zasobowe. Taka struktura umożliwia precyzyjne, selektywne skalowanie wyłącznie tych komponentów, które w danym momencie stanowią wąskie gardło systemu.

\subsection{Wzorce skalowania: pionowe i poziome}
\label{subsec:wzorce_skalowania}
W inżynierii systemów chmurowych wyróżnia się dwa fundamentalne kierunki adaptacji zasobowej:
\begin{itemize}
    \item \textbf{Skalowanie pionowe (ang. \textit{scaling up}):} polega na zwiększaniu zasobów sprzętowych (dodawanie rdzeni CPU, rozbudowa pamięci RAM) istniejącego węzła lub kontenera. Metoda ta posiada krytyczne ograniczenia: wiąże się z fizycznymi barierami sprzętowymi maszyn wirtualnych oraz najczęściej wymaga restartu procesu aplikacji, co generuje niedopuszczalny w systemach SaaS czas przestoju (ang. \textit{downtime}).
    \item \textbf{Skalowanie poziome (ang. \textit{scaling out}):} polega na dynamicznym dodawaniu kolejnych równoległych instancji aplikacyjnych (Podów/kontenerów) współpracujących za warstwą równoważenia obciążenia (ang. \textit{Load Balancer}). Jest to wzorzec natywny dla chmury, pozwalający na niemal nieograniczoną rozbudowę wydajnościową systemu bez przerywania obsługi żądań użytkowników końcowych.
\end{itemize}

\section{Platformy orkiestracji kontenerów}
\label{sec:platformy_orkiestracji}

\subsection{Kubernetes -- architektura i kluczowe komponenty}
\label{subsec:k8s_architektura_wstep}
Kubernetes (K8s) jest otwartoźródłową platformą służącą do automatyzacji wdrażania, skalowania i zarządzania aplikacjami skonteneryzowanymi. System opiera się na architekturze klastrowej o podziale na płaszczyznę kontrolną (\textit{Control Plane}) oraz węzły wykonawcze (\textit{Worker Nodes}). Płaszczyzna kontrolna odpowiada za nadrzędne sterowanie klastrem: serwer API (\texttt{kube-apiserver}) stanowi centralny punkt komunikacyjny, \texttt{etcd} przechowuje stan klastra jako rozproszona baza klucz-wartość, a \texttt{kube-scheduler} przypisuje nowo utworzone kontenery do fizycznych węzłów.

Na węzłach wykonawczych kluczową rolę odgrywa demon \texttt{kubelet}, który nadzoruje stan kontenerów i komunikuje się z serwerem API, oraz \texttt{kube-proxy}, odpowiedzialny za niskopoziomowe mapowanie ruchu sieciowego i routing. Podstawową jednostką orkiestracji w Kubernetesie jest Pod -- logiczna grupa jednego lub wielu kontenerów współdzielących przestrzeń sieciową (adres IP) oraz wolumeny dyskowe.

\subsection{Ewolucja środowisk badawczych: od symulacji lokalnych do chmury publicznej}
\label{subsec:ewolucja_srodowisk}
W fazie projektowania systemów rozproszonych powszechną praktyką jest stosowanie lekkich, lokalnych dystrybucji orkiestratora, takich jak \textit{Minikube} lub \textit{k3s} (zoptymalizowana dystrybucja CNCF pozbawiona nadmiarowych sterowników dostawców chmurowych). Środowiska te doskonale sprawdzają się w procesach ciągłej integracji (CI/CD) oraz podczas weryfikacji poprawności logicznej manifestów aplikacyjnych.

Jednakże, z punktu widzenia precyzyjnej metrologii i testów wydajnościowych systemów klasy SaaS, symulacje lokalne wprowadzają szereg zakłóceń uniemożliwiających prawidłową ocenę fizyki zjawisk chmurowych. Współdzielenie jądra systemu operacyjnego maszyny hosta, brak rzeczywistej izolacji sprzętowej, uproszczona warstwa sieciowa oparta na mostkach wewnątrz-systemowych (w przeciwieństwie do routingu chmurowego VPC) oraz natychmiastowa, nierealistyczna alokacja zasobów dyskowych sprawiają, że kluczowe anomalie dynamiczne pozostają niewykrywalne. W związku z tym, w celu zapewnienia wysokiej wierności pomiarowej, jako środowisko badawcze w niniejszej pracy zaimplementowano pełnowymiarowy, zarządzany klaster \textit{Google Kubernetes Engine (GKE)}, eliminując uproszczenia lokalne na rzecz produkcyjnych mechanizmów chmury publicznej.

\section{Struktura dokumentu}
\label{sec:struktura_dokumentu}
Praca została podzielona na siedem powiązanych logicznie rozdziałów:
\begin{*}{\label{item:struktura}}
    \item \textbf{Rozdział 1} wprowadza w tematykę pracy, definiuje hipotezę badawczą oraz cel i zakres realizowanych eksperymentów.
    \item \textbf{Rozdział 2} zawiera szczegółowy przegląd technologii chmurowych, strukturę matematyczną algorytmu HPA oraz opis aparatury pomiarowej.
    \item \textbf{Rozdział 3} przedstawia projekt architektury systemu SaaS, charakterystykę profili obciążeniowych mikroserwisów oraz strukturę macierzy badawczej.
    \item \textbf{Rozdział 4} dokumentuje proces implementacji środowiska w chmurze GCP, konfigurację skryptów k6 z pominięciem DNS oraz deklaratywne manifesty orkiestracji.
    \item \textbf{Rozdział 5} stanowi rdzeń empiryczny pracy -- zawiera zbiorcze zestawienie wyników pomiarowych oraz drobiazgową analizę zidentyfikowanych anomalii chmurowych.
    \item \textbf{Rozdział 6} dedykowany jest zagadnieniom izolacji i skalowania w architekturach wielodostępnych (\textit{Multi-Tenancy}).
    \item \textbf{Rozdział 7} podsumowuje uzyskane wyniki, weryfikuje hipotezę oraz formułuje rekomendacje inżynierskie dla projektantów chmurowych systemów SaaS.
\end{*}